\noindent
This report provides a comprehensive, first-principles treatment of the optical
characterisation of a three-period [TiO$_2$/Ni]$\times$3 hyperbolic metamaterial
(HMM) deposited on a borosilicate glass (SiO$_2$) substrate. Starting from
Maxwell's equations, we derive the uniaxial electromagnetic dispersion relation,
effective-medium theory (EMT), and the complete transfer matrix method (TMM) for
anisotropic multilayer stacks, applied step by step to the FIB/TEM-measured
per-layer thicknesses of the fabricated sample. We document systematically why
angle-resolved reflectance and transmittance measurements are fundamentally
insufficient to recover the out-of-plane permittivity $\eperp$ at sub-wavelength
film thicknesses---a consequence of the evanescent suppression of the $E_z$ field
component across a \mbox{44\,nm} stack---and how spectroscopic ellipsometry resolves
this through phase-sensitive multi-angle inversion. The extracted permittivity
tensor $(\epar(\lambda),\,\eperp(\lambda))$ is used to reconstruct the
isofrequency contour (IFC) via the Jackson et al.\ (2024) lossy complex-wavevector
method. The central finding is that both $\Re(\epar)$ and $\Re(\eperp)$ remain
positive throughout $550$--$1000$\,nm, producing a closed elliptical IFC with no
hyperbolic transition in the measured spectral window.
